% ============ 第1章 §1.6 锐极大函数与Hardy-Littlewood极大函数的关系(原书 pp.44-51 = PDF pp.56-63) ============
\section{锐极大函数与 Hardy--Littlewood 极大函数的关系}\label{sec:1.6}
\index{good $\lambda$-不等式}
\index{Whitney 型覆盖定理}
\index{二进区间}


\begin{definition}\label{def:1.6.1}
形如 $[\ell 2^{k}, (\ell + 1)2^{k})$(其中 $k, \ell \in \mathbb{Z}$)的区间称为二进区间.令 $\mathbb{D}_{k} = \{[\ell 2^{k}, (\ell + 1)2^{k})\}_{\ell \in \mathbb{Z}}$,并用 $\mathbb{D}$ 表示全体二进区间.
\end{definition}

对二进区间,成立如下基本性质.

\begin{lemma}\label{lem:1.6.2}
(i) 对每个 $x \in \mathbb{R}$ 与每个 $k \in \mathbb{Z}$,包含 $x$ 的二进区间 $Q \in \mathbb{D}_{k}$ 恰有一个.

(ii) 若 $Q, Q' \in \mathbb{D}$,则 $Q \subset Q'$,$Q' \subset Q$ 或 $Q \cap Q' = \emptyset$ 三者恰有一个发生.

(iii) 对每个 $Q \in \mathbb{D}_{k}$,包含 $Q$ 的 $Q' \in \mathbb{D}_{k+1}$ 恰有一个.
\end{lemma}

由这个几乎自明的事实可得出下面的事实.

\begin{lemma}\label{lem:1.6.3}
设 $Q_{j} \in \mathbb{D}$($j = 1, 2, \ldots$),且在 $\{Q_{j}\}_{j \in \mathbb{N}}$ 中不存在无限扩张的子区间列.则可以从 $\{Q_{j}\}_{j \in \mathbb{N}}$ 中选出子区间列 $\{Q_{j_{\ell}}\}_{\ell}$(有限个或无限个),使其满足
\begin{enumerate}[(i)]
\item $\bigcup_{j \in \mathbb{N}}Q_{j} = \bigcup_{\ell}Q_{j_{\ell}}$,
\item $Q_{j_{\ell}} \cap Q_{j_{k}} = \emptyset$($\ell \neq k$),
\item $Q_{j_{\ell}}$ 在 $\{Q_{j}\}_{j \in \mathbb{N}}$ 中是极大的,即 $\{Q_{j}\}_{j \in \mathbb{N}}$ 中不存在真包含 $Q_{j_{\ell}}$ 的区间.
\end{enumerate}
\end{lemma}

\begin{proof}
由引理~\ref{lem:1.6.2}(ii),(iii),对每个 $Q_{j}$,包含 $Q_{j}$ 的区间 $Q_{k}$ 的全体按包含关系排成一列.由假设,这一列是有限列,故最宽的区间唯一确定,它在包含关系下是极大的.把在包含关系下极大的这些 $Q_{k}$ 的全体记为 $\{Q_{j_{\ell}}\}$,则
\begin{equation*}
\bigcup_{j \in \mathbb{N}}Q_{j} = \bigcup_{\ell}Q_{j_{\ell}},
\end{equation*}
且由极大性,
\begin{equation*}
Q_{j_{\ell}} \cap Q_{j_{k}} = \emptyset \quad (\ell \neq k),
\end{equation*}
引理的结论全部得到满足.
\end{proof}

利用二进区间的这些基本性质,可以证明如下的 Whitney 型覆盖定理.

\begin{theorem}[Whitney 型覆盖定理]\label{thm:1.6.4}
设 $\Omega$ 是 $\mathbb{R}$ 的非空真开子集.则存在闭区间列 $\{Q_{j}\}$(有限个或无限个)满足:
\begin{enumerate}[(i)]
\item $\Omega = \bigcup_{j}Q_{j}$ 且 $Q_{j}^{\circ} \cap Q_{k}^{\circ} = \emptyset$($j \neq k$);
\item $\ell(Q_{j}) \leq \mathrm{dist}(Q_{j}, \Omega^{c}) \leq 4\ell(Q_{j})$;
\item 若 $Q_{j}$ 与 $Q_{k}$ 相接,则
\begin{equation*}
\frac{1}{4} \leq \frac{\ell(Q_{j})}{\ell(Q_{k})} \leq 4.
\end{equation*}
\end{enumerate}
\end{theorem}

\begin{proof}
对 $k \in \mathbb{Z}$,令
\begin{equation*}
\Omega_{k} = \{x \in \Omega;\ 2\cdot 2^{-k} < \mathrm{dist}(x, \Omega^{c}) \leq 4\cdot 2^{-k}\},
\end{equation*}
则
\begin{equation*}
\Omega = \bigcup_{k \in \mathbb{Z}}\Omega_{k}.
\end{equation*}
设 $E_{k}$ 是满足 $Q \cap \Omega_{k} \neq \emptyset$ 且 $Q \in \mathbb{D}_{-k}$ 的二进区间的全体,并令 $E = \bigcup_{k \in \mathbb{Z}}E_{k}$.

若 $Q \in E$,则对某个 $k \in \mathbb{Z}$ 有 $Q \in E_{k}$,故对这个 $k$,存在 $x \in Q \cap \Omega_{k}$,且 $\ell(Q) = 2^{-k}$.由 $x \in \Omega_{k}$,$2\cdot 2^{-k} < \mathrm{dist}(x, \Omega^{c}) \leq 4\cdot 2^{-k}$,所以
\begin{equation*}
\ell(Q) < \mathrm{dist}(x, \Omega^{c}) - \ell(Q) \leq \mathrm{dist}(Q, \Omega^{c}) \leq \mathrm{dist}(x, \Omega^{c}) \leq 4\cdot 2^{-k} = 4\ell(Q),
\end{equation*}
故 $Q$ 满足条件 (ii).又,$\mathrm{dist}(Q, \Omega^{c}) \geq \ell(Q) > 0$,故 $Q \subset \Omega$.因此 $\bigcup_{Q \in E}Q \subset \Omega$.

反之,若 $x \in \Omega$,则存在 $k \in \mathbb{Z}$ 使得 $x \in \Omega_{k}$.进一步,由于 $\bigcup_{Q \in \mathbb{D}_{-k}}Q = \mathbb{R}$,存在 $Q \in \mathbb{D}_{-k}$ 使得 $x \in Q$.由 $E_{k}$ 的定义,这个 $Q$ 属于 $E_{k}$,从而属于 $E$.因此,$\Omega \subset \bigcup_{Q \in E}Q$,得
\begin{equation*}
\Omega = \bigcup_{Q \in E}Q.
\end{equation*}

又,若 $Q, Q' \in E$,则分别有
\begin{equation*}
\ell(Q) < \mathrm{dist}(Q, \Omega^{c}) \leq 4\ell(Q), \quad \ell(Q') < \mathrm{dist}(Q', \Omega^{c}) \leq 4\ell(Q'),
\end{equation*}
进一步设 $Q \subset Q'$,注意到 $\mathrm{dist}(Q, \Omega^{c}) \geq \mathrm{dist}(Q', \Omega^{c})$,则
\begin{equation*}
\ell(Q') \leq 4\ell(Q).
\end{equation*}
因此,$E$ 中包含给定 $Q \in E$ 的二进区间至多 2 个.也就是说,$E$ 中的二进区间满足引理~\ref{lem:1.6.3} 的假设.于是,设由引理~\ref{lem:1.6.3} 从 $E$ 中得到的极大二进区间列为 $\{Q_{j}\}$,则
\begin{equation*}
\Omega = \bigcup_{j}Q_{j} \quad \text{且} \quad Q_{j} \cap Q_{k} = \emptyset \quad (j \neq k).
\end{equation*}

现在来确认 (iii).若 $Q_{j}$ 与 $Q_{k}$ 相接,则
\begin{equation*}
\begin{aligned}
\ell(Q_{j}) &\leq \mathrm{dist}(Q_{j}, \Omega^{c}) \leq \mathrm{diam}(Q_{k}) + \mathrm{dist}(Q_{k}, \Omega^{c}) \\
&\leq \ell(Q_{k}) + 4\ell(Q_{k}) \leq 5\ell(Q_{k}).
\end{aligned}
\end{equation*}
因此 $\ell(Q_{j})/\ell(Q_{k}) \leq 5$.但是,$Q_{j}$,$Q_{k}$ 都是二进区间,故 $\ell(Q_{j})/\ell(Q_{k})$ 具有 $2^{m}$($m \in \mathbb{Z}$)的形状,所以 $\ell(Q_{j})/\ell(Q_{k}) \leq 4$.交换 $Q_{j}$ 与 $Q_{k}$ 的角色,得 $\ell(Q_{k})/\ell(Q_{j}) \leq 4$,从而 $1/4 \leq \ell(Q_{j})/\ell(Q_{k}) \leq 4$.若把 $Q_{j}$ 换成加上其右端点的闭区间(仍记作 $Q_{j}$),则 $\{Q_{j}\}$ 满足全部所需性质.
\end{proof}

利用这个覆盖定理,可证明给出 Hardy--Littlewood 极大函数与锐极大函数关系的如下 ``good $\lambda$-不等式''.

\begin{lemma}\label{lem:1.6.5}
若 $f \in L_{\mathrm{loc}}^{1}(\mathbb{R})$,则对任意 $\lambda > 0$,$\gamma > 0$,
\begin{equation*}
\left|\left\{x \in \mathbb{R};\ M_{1}f(x) > 2\lambda,\ f^{\#}(x) < \gamma\lambda\right\}\right| \leq 96\gamma\left|\left\{x \in \mathbb{R};\ M_{1}f(x) > \lambda\right\}\right|
\end{equation*}
成立.
\end{lemma}

\begin{proof}
以下,将 $M_{1}f$ 简记为 $Mf$.设 $\lambda > 0$,$\gamma > 0$,令 $E_{\lambda} = \{x \in \mathbb{R};\ Mf(x) > \lambda\}$.若 $E_{\lambda} = \mathbb{R}$,则引理的结论正确,故以下设 $E_{\lambda} \neq \mathbb{R}$.$E_{\lambda}$ 是开集.事实上,若 $x_{0} \in E_{\lambda}$,由上确界的定义,存在包含 $x_{0}$ 的区间 $Q_{0}$ 使得
\begin{equation*}
\frac{1}{|Q_{0}|}\int_{Q_{0}}|f(x)|\dd x > \lambda.
\end{equation*}
由积分的性质,可取一个包含 $Q_{0}$ 的开区间 $Q_{1}$ 仍满足上述性质.若 $x \in Q_{1}$,则由 $\frac{1}{|Q_{1}|}\int_{Q_{1}}|f(x)|\dd x > \lambda$ 得 $Mf(x) > \lambda$,故 $Q_{1} \subset E_{\lambda}$.因此 $E_{\lambda}$ 是开集.又因 $E_{\lambda} \neq \mathbb{R}$,$E_{\lambda}$ 满足定理~\ref{thm:1.6.4} 的假设,故可取如下区间列 $\{Q_{j}\}$:
\begin{equation*}
(\mathrm{i})\ E_{\lambda} = \bigcup_{j}\mathring{Q}_{j} \quad \text{且} \quad Q_{j}^{\circ} \cap Q_{k}^{\circ} = \emptyset \quad (j \neq k).
\end{equation*}
\begin{equation*}
\text{(ii)}\ \ell(Q_{j}) \leq \mathrm{dist}(Q_{j}, E_{\lambda}^{c}) \leq 4\ell(Q_{j}).
\end{equation*}
固定其中任意一个 $Q_{j}$,为简单起见记作 $Q$.

设 $\tilde{Q} = 12Q$.由于 $\mathrm{dist}(Q, E_{\lambda}^{c}) \leq 4\ell(Q) = \frac{1}{3}\ell(\tilde{Q})$,故 $\tilde{Q} \cap E_{\lambda}^{c} \neq \emptyset$,由 $E_{\lambda}$ 的定义,
\begin{equation*}
\frac{1}{|\tilde{Q}|}\int_{\tilde{Q}}|f(x)|\dd x \leq \lambda.
\end{equation*}

其次,若 $x \in Q$ 且 $Mf(x) > 2\lambda$,则存在包含 $x$ 的区间 $R$ 使得 $|R|^{-1}\int_{R}|f(y)|\dd y > 2\lambda$.于是,设 $R'$ 是以 $x$ 为中心,长度为 $2\ell(R)$ 的区间,则
\begin{equation*}
\frac{1}{|R'|}\int_{R'}|f(y)|\dd y \geq \frac{|R|}{|R'|}\cdot\frac{1}{|R|}\int_{R}|f(y)|\dd y > \lambda.
\end{equation*}
若 $R' \cap E_{\lambda}^{c} \neq \emptyset$,则 $|R'|^{-1}\int_{R'}|f(y)|\dd y \leq \lambda$,与上面的不等式矛盾,故 $R' \cap E_{\lambda}^{c} = \emptyset$.又由 $\mathrm{dist}(Q, E_{\lambda}^{c}) \leq 4\ell(Q)$,
\begin{equation*}
\mathrm{dist}(x, E_{\lambda}^{c}) \leq \ell(Q) + \mathrm{dist}(Q, E_{\lambda}^{c}) \leq 5\ell(Q).
\end{equation*}
因此 $R' \subset Q(x, 11\ell(Q))$,从而 $R \subset Q(x, 11\ell(Q))$.用 $x_{Q}$ 表示 $Q$ 的中心,则 $|x - x_{Q}| \leq \ell(Q)/2$,故 $R \subset Q(x, 11\ell(Q)) \subset Q(x_{Q}, 12\ell(Q)) = \tilde{Q}$.因此 $Q \subset \tilde{Q}$ 且 $R \subset \tilde{Q}$.于是,
\begin{equation*}
\frac{1}{|R|}\int_{R}|f(y)\chi_{\tilde{Q}}(y)|\dd y = \frac{1}{|R|}\int_{R}|f(y)|\dd y > 2\lambda
\end{equation*}
所以
\begin{align*}
\frac{1}{|R|}\int_{R}|f(y) - f_{\tilde{Q}}|\chi_{\tilde{Q}}(y)\dd y&= \frac{1}{|R|}\int_{R}|f(y) - f_{\tilde{Q}}|\dd y \geq \frac{1}{|R|}\int_{R}|f(y)|\dd y - |f_{\tilde{Q}}|  \\
\qquad &\geq \frac{1}{|R|}\int_{R}|f(y)|\dd y - \frac{1}{|\tilde{Q}|}\int_{\tilde{Q}}|f(y)|\dd y > 2\lambda - \lambda = \lambda.
\end{align*}
这表明
\begin{equation*}
M\bigl((f - f_{\tilde{Q}})\chi_{\tilde{Q}}\bigr)(x) > \lambda.
\end{equation*}

对 $M\bigl((f - f_{\tilde{Q}})\chi_{\tilde{Q}}\bigr)$ 使用定理~\ref{thm:1.5.2},
\begin{align*}
|\{x \in \mathbb{R};\ M\bigl((f - f_{\tilde{Q}})\chi_{\tilde{Q}}\bigr)(x) > \lambda\}|&\leq \frac{8}{\lambda}\int_{\tilde{Q}}|f(y) - f_{\tilde{Q}}|\dd y \leq \frac{8|\tilde{Q}|}{\lambda}\cdot\frac{1}{|\tilde{Q}|}\int_{\tilde{Q}}|f(y) - f_{\tilde{Q}}|\dd y  \\
&\leq \frac{8\cdot 12|Q|}{\lambda}\inf_{x \in \tilde{Q}}f^{\#}(x) \leq \frac{96|Q|}{\lambda}\inf_{x \in Q}f^{\#}(x).
\end{align*}
因此,若 $\inf_{x \in Q}f^{\#}(x) < \gamma\lambda$,则
\begin{align*}
& |\{x \in Q;\ Mf(x) > 2\lambda,\ f^{\#}(x) < \gamma\lambda\}| \\
& \qquad \leq |\{x \in Q;\ M\bigl((f - f_{\tilde{Q}})\chi_{\tilde{Q}}\bigr)(x) > \lambda,\ f^{\#}(x) < \gamma\lambda\}| \\
& \qquad \leq \frac{96|Q|}{\lambda}\inf_{x \in Q}f^{\#}(x) = 96\gamma|Q|.
\end{align*}
又,若 $\inf_{x \in Q}f^{\#}(x) \geq \gamma\lambda$,则上式左端内的集合为空,故上式成立.无论哪种情形,
\begin{equation}\tag{1.6.1}
\left|\{x \in Q;\ Mf(x) > 2\lambda,\ f^{\#}(x) < \gamma\lambda\}\right| \leq 96\gamma|Q|
\label{eq:1.6.1}\end{equation}
得证.

把对 $Q_{j}$ 的估计式~\eqref{eq:1.6.1} 对 $Q_{j}$ 求和,由于 $Q_{j}$ 的内部两两不相交,
\begin{align*}
& |\{x \in \mathbb{R};\ Mf(x) > 2\lambda,\ f^{\#}(x) < \gamma\lambda\}| \\
& \qquad = |\{x \in E_{\lambda};\ Mf(x) > 2\lambda,\ f^{\#}(x) < \gamma\lambda\}| \\
& \qquad \leq 96\gamma|E_{\lambda}| = 96\gamma|\{x \in \mathbb{R};\ Mf(x) > \lambda\}|.
\end{align*}
\end{proof}

下面给出联系 ``good $\lambda$-不等式'' 与 $L^{p}$ 不等式的引理.

\begin{lemma}\label{lem:1.6.6}
设 $0 \leq w(x) \in L_{\mathrm{loc}}^{1}(\mathbb{R})$,$0 < p < \infty$.设 $f$,$g$ 是可测函数,$|f(x)|^{p}w(x) \in L^{1}(\mathbb{R})$,且存在满足以下两个条件的 $\varepsilon > 0$,$\gamma > 0$:
\begin{enumerate}[(i)]
\item $(1 + \varepsilon)^{p}\gamma < 1$,
\item $w(\{x \in \mathbb{R};\ |f(x)| > (1 + \varepsilon)\lambda,\ |g(x)| \leq \lambda\}) \leq \gamma w(\{x \in \mathbb{R};\ |f(x)| > \lambda\})$($\lambda > 0$).
\end{enumerate}
其中,$w(E) = \int_{E}w(x)\dd x$.

则存在常数 $C(p, \varepsilon, \gamma)$ 使得
\begin{equation*}
\| f \|_{L^{p}(w)} \leq C(p, \varepsilon, \gamma)\| g \|_{L^{p}(w)}.
\end{equation*}
其中,$\| f \|_{L^{p}(w)} = \left(\int_{\mathbb{R}}|f(x)|^{p}w(x)\dd x\right)^{1/p}$.
\end{lemma}

\begin{proof}
由假设 (ii),
\begin{align*}
w(\{x \in \mathbb{R};\ |f(x)| > (1 + \varepsilon)\lambda\})&\leq w(\{x \in \mathbb{R};\ |g(x)| > \lambda\}) + \gamma w(\{x \in \mathbb{R};\ |f(x)| > \lambda\})
 \end{align*}
所以
\begin{align*}
\| f \|_{L^{p}(w)}^{p} &= \int_{\mathbb{R}}|f(x)|^{p}w(x)\dd x \\
\qquad &= \int_{0}^{\infty}p\lambda^{p-1}w(\{x \in \mathbb{R};\ |f(x)| > \lambda\})\dd \lambda \\
\qquad &= (1 + \varepsilon)^{p}\int_{0}^{\infty}p\lambda^{p-1}w(\{x \in \mathbb{R};\ |f(x)| > (1 + \varepsilon)\lambda\})\dd \lambda \\
\qquad &\leq (1 + \varepsilon)^{p}\int_{0}^{\infty}p\lambda^{p-1}w(\{x \in \mathbb{R};\ |g(x)| > \lambda\})\dd \lambda \\
&\qquad\quad + \gamma(1 + \varepsilon)^{p}\int_{0}^{\infty}p\lambda^{p-1}w(\{x \in \mathbb{R};\ |f(x)| > \lambda\})\dd \lambda \\
\qquad &= (1 + \varepsilon)^{p}\| g \|_{L^{p}(w)}^{p} + \gamma(1 + \varepsilon)^{p}\| f \|_{L^{p}(w)}^{p}.
\end{align*}
因此,由 $|f(x)|^{p}w(x) \in L^{1}(\mathbb{R})$ 得 $\| f \|_{L^{p}(w)} < \infty$,从两边减去 $\gamma(1 + \varepsilon)^{p}\| f \|_{L^{p}(w)}^{p}$,得
\begin{equation*}
\left(1 - \gamma(1 + \varepsilon)^{p}\right)\| f \|_{L^{p}(w)}^{p} \leq (1 + \varepsilon)^{p}\| g \|_{L^{p}(w)}^{p},
\end{equation*}
由条件 $(1 + \varepsilon)^{p}\gamma < 1$,
\begin{equation*}
\| f \|_{L^{p}(w)} \leq \frac{(1 + \varepsilon)}{(1 - \gamma(1 + \varepsilon)^{p})^{1/p}}\| g \|_{L^{p}(w)}.
\end{equation*}
\end{proof}

\begin{note}\label{note:1.6.7}
条件 $f \in L^{p}(w)$ 也可以用 $\inf(1, |f|) \in L^{p}(w)$ 代替.进一步,更一般地,也可以用存在 $\alpha > 0$ 与 $C > 0$ 使得
\begin{equation*}
w(\{x \in \mathbb{R};\ |f(x)| > \lambda\}) \leq \frac{C}{\lambda^{\alpha}} \quad (0 < \lambda < 1)
\end{equation*}
这一条件来代替(宫地,薮田,On good $\lambda$-inequality, Bull.\ Fac.\ Sci.\ Ibaraki Univ.\ 16 (1984), 1--11).
\end{note}

利用关于 Hardy--Littlewood 极大函数与锐极大函数的 ``good $\lambda$-不等式'' 以及上面的引理,可证明给出 Hardy--Littlewood 极大函数与锐极大函数 $L^{p}$ 关系的如下定理.

\begin{theorem}\label{thm:1.6.8}
设 $1 \leq p < \infty$.若 $f \in L_{\mathrm{loc}}^{1}(\mathbb{R})$ 且 $Mf \in L^{p}(\mathbb{R})$,则存在与 $f$ 无关的常数 $C_{p} > 0$ 使得
\begin{equation*}
\| Mf \|_{L^{p}} \leq C_{p}\| f^{\#} \|_{L^{p}}.
\end{equation*}
\end{theorem}

\begin{proof}
由引理~\ref{lem:1.6.5},
\begin{equation*}
|\{x \in \mathbb{R};\ Mf(x) > 2\lambda,\ \frac{f^{\#}(x)}{\gamma} < \lambda\}| \leq 96\gamma|\{x \in \mathbb{R};\ Mf(x) > \lambda\}| \quad (\lambda > 0).
\end{equation*}
于是固定一个 $\gamma$ 使得 $2^{p} \times 96\gamma < 1$.则,由引理~\ref{lem:1.6.6},
\begin{align*}
\| Mf \|_{L^{p}} & \leq \frac{2}{(1 - 2^{p} \times 96\gamma)^{1/p}}\left\| \frac{f^{\#}(x)}{\gamma} \right\|_{L^{p}} \\
& \leq \frac{2}{\gamma(1 - 2^{p} \times 96\gamma)^{1/p}}\| f^{\#} \|_{L^{p}}.
\end{align*}
\end{proof}

\begin{note}\label{note:1.6.9}
条件 $Mf \in L^{p}(\mathbb{R})$ 也可以用 $\inf(1, |f|) \in L^{p}(\mathbb{R})$ 代替.进一步,更一般地,也可以用
\begin{equation*}
|\{x \in \mathbb{R};\ |f(x)| > \lambda\}| < +\infty \quad (0 < \lambda < 1)
\end{equation*}
这一条件来代替(参见注意~\ref{note:1.6.7} 的文献).
\end{note}
